\chapter{抓取规划、抓取质量与非抓持操作}
\label{ch:grasp-planning}

% ════════════════════════════════════════════════════════════════
%  基于正典重述（教材化、记号归一、非翻译非抄录）：
%   源 MLS Murray-Li-Sastry Ch5《多指手运动学》§3-4,§6（抓取静力学/力封闭/
%       抓取规划/滚动接触运动学·Montana 方程）+ Modern Robotics Lynch-Park
%       §12.1.7（形封闭与 LP 判据/质量度量）、§12.2.3（力封闭/三接触约化）、
%       §12.3（拟静态推动·meter-stick·动态抓取·峰嵌）。
%  承 \cref{ch:grasp-closure}（接触建模/摩擦锥/形封闭力封闭基础），本章只做
%       「规划与度量」与「非抓持/滚动」，与其互引不重复。
%  记号归一（与本书全局及 cooperative_manipulation 一致）：
%   向量粗体小写、矩阵粗体大写；R in SO(3)、T in SE(3)（preamble 宏 \SO \SE \so \se）；
%   抓取映射 w=Gf、点接触 G_i=[I; r̂_i]（\cref{eq:cm-Gi}）；wrench=力旋量、twist=运动旋量；
%   se(3) 旋量 ξ=[rho;φ]（平移在前）；tau=Jᵀf；右扰动主线。
% ════════════════════════════════════════════════════════════════

% —— 本章局部数学算子（章文件 \input 入正文，故用 \providecommand+\operatorname 兜底，
%    与既有 preamble/兄弟章定义不冲突）——
\providecommand{\rank}{\operatorname{rank}}
\providecommand{\pos}{\operatorname{pos}}      % 正张成（锥包）
\providecommand{\conv}{\operatorname{conv}}    % 凸包
\providecommand{\spanop}{\operatorname{span}}  % 线性张成
\providecommand{\interior}{\operatorname{int}} % 内部（避开 plain-TeX 内置 \intop）
\providecommand{\trace}{\operatorname{tr}}
\providecommand{\diagop}{\operatorname{diag}}
% 反对称（叉积）矩阵：位置向量 r 的 \widehat{r}（与 \cref{ch:lie} 的 (·)^\wedge 同义）
\providecommand{\skewm}[1]{\widehat{#1}}

\begin{practice}[前置自测]
开始之前，先试答下列问题。若已能流畅作答，可只速读本章的判据卡与小结；若卡住，正是本章要为你解决的。
\begin{enumerate}
  \item 你已知（\cref{ch:grasp-closure}）平面物体形封闭至少需 $4$ 个点接触、空间需 $7$ 个。但若每个接触\emph{带摩擦}，所需手指数为何骤减到 $2\sim3$？“形封闭”与“力封闭”到底差在哪一项物理假设？
  \item 同一个物体有无数种力封闭抓取，凭什么说“这个抓取比那个好”？请给出一个\emph{可计算的标量}来排序两个抓取的优劣。
  \item 两根手指捏住一块砖，连接两接触点的直线必须落在两个摩擦锥内——这个“对踵（antipodal）”条件为什么\emph{恰好}等价于二指力封闭？
  \item 一根手指\emph{不夹}、只\emph{推}地上的盒子，盒子会怎么转？决定它绕哪点转的，是质量分布还是接触摩擦的几何？
  \item 软指尖在物体曲面上\emph{滚动}时，接触点在两个曲面上各自爬行。给定两曲面的曲率，接触点坐标随相对运动如何演化？为什么平面上滚一个球，接触点轨迹只由球半径决定？
\end{enumerate}
\noindent\emph{答案线索}：1、3 承 \cref{ch:grasp-closure} 的摩擦锥与 \cref{def:friction-cone}；2 即本章 \cref{sec:gp-quality} 的抓取质量度量；4 是 \cref{sec:gp-nonprehensile} 的拟静态推动；5 是 \cref{sec:gp-rolling} 的 Montana 接触方程。
\end{practice}

\section*{本章目标}
学完本章，你应当能够：
\begin{enumerate}
  \item \textbf{度量}一个力封闭抓取的“好坏”：写出抓取力旋量集（grasp wrench set），用\emph{最大内接球半径}（largest-minimum-resisted-wrench）与 Ferrari--Canny $\varepsilon$ 度量给抓取打分，并理解面向任务的度量如何用任务力旋量空间（task wrench space）替换“各向同性球”；
  \item \textbf{判定与构造}力封闭抓取：陈述并理解二指对踵抓取定理（平面充要、空间软指推广），用 Caratheodory/Steinitz 定理给出力封闭所需接触数的上下界，复述“把不满足约束的接触拉回到力封闭区域”的构造算法思想；
  \item \textbf{分析}非抓持/准静态操作：用拟静态力平衡（quasistatic force balance）与摩擦极限面（limit surface）分析推动（pushing）的运动学，解释 meter-stick 翻转与峰嵌（wedging），并把握动态抓取（dynamic grasp）的惯性力机理；
  \item \textbf{推导}滚动接触运动学：从曲面第一/二基本形式与曲率张量出发，在归一化 Gauss 标架下写出 Montana 接触方程，求接触坐标 $(\boldsymbol{\alpha}_f,\boldsymbol{\alpha}_o,\psi)$ 随相对运动旋量的演化，并算通“球在平面上滚”的小例；
  \item \textbf{衔接}本章度量与构造到协同操作（\cref{ch:mr-coopforce}）的抓取矩阵 $\mathbf{G}$、闭链双臂规划（\cref{ch:arm-dualarm-planning}）的约束流形。
\end{enumerate}

\subsection*{本章知识导航}
本章是一条“\emph{从能不能抓（\cref{ch:grasp-closure}）走向抓得好不好、怎么抓、不抓怎么动}”的主线。前一章（\cref{ch:grasp-closure}）已奠定接触建模、摩擦锥、形封闭/力封闭的\emph{二值判据}（封闭/不封闭）；本章把二值判据升级为\emph{连续度量}（抓取质量），给出\emph{构造算法}（对踵/枚举），再越出“抓持”进入\emph{非抓持操作}（推、倒、滚）。

\begin{center}
\small
\begin{tikzpicture}[
  node distance=5mm and 8mm, >=Stealth,
  blk/.style={draw, rounded corners, align=center, font=\small, inner sep=3pt, fill=main!6, draw=main!60!black},
  grp/.style={draw, rounded corners, align=center, font=\small, inner sep=3pt, fill=second!6, draw=second!60!black},
  ln/.style={->, thick, gray!60!black}]
\node[blk] (cl) {封闭判据\\ \cref{ch:grasp-closure}};
\node[blk, right=of cl] (q) {抓取质量\\ $\varepsilon$ 度量};
\node[blk, right=of q] (plan) {抓取规划\\ 对踵·枚举};
\node[grp, below=8mm of cl] (push) {拟静态推动\\ 极限面};
\node[grp, right=of push] (flip) {倾倒/翻转\\ 动态抓取};
\node[grp, right=of flip] (roll) {滚动接触\\ Montana 方程};
\node[blk, below=8mm of push, fill=third!8, draw=third!60!black] (coop) {抓取矩阵 $\mathbf{G}$\\ \cref{ch:mr-coopforce}};
\node[blk, right=of coop, fill=third!8, draw=third!60!black] (dual) {约束流形\\ \cref{ch:arm-dualarm-planning}};
\draw[ln] (cl)--(q); \draw[ln] (q)--(plan);
\draw[ln] (cl.south) -- ++(0,-4mm) -| (push.north);
\draw[ln] (push)--(flip); \draw[ln] (flip)--(roll);
\draw[ln] (q.south) -- ++(0,-22mm) -| (coop.north);
\draw[ln] (plan.south) -- ++(0,-26mm) -| (dual.north);
\end{tikzpicture}
\end{center}

\begin{insight}[本章一句话]
\cref{ch:grasp-closure} 回答“接触\emph{能否}锁住物体”；本章回答“锁得\emph{多牢}（质量度量）、\emph{怎样}选接触点锁住（规划）、以及\emph{不锁}时如何驱动物体（非抓持）”。三件事共用同一套工具——把接触力的可行集写成\emph{凸锥/凸多胞体}，再在其上做几何（内接球、凸包、力平衡）。
\end{insight}

\paragraph{承上：抓取映射 $\mathbf{G}$ 是全章的公共语言。}本章自始至终复用\cref{ch:mr-coopforce}建立的抓取映射 $\mathbf{w}=\mathbf{G}\mathbf{f}$（\cref{eq:cm-grasp}）：$N$ 个接触施加的接触力 $\mathbf{f}=[\mathbf{f}_1^\top,\dots,\mathbf{f}_N^\top]^\top$ 经抓取矩阵 $\mathbf{G}$ 合成作用于物体质心的力旋量 $\mathbf{w}=[\mathbf{F}^\top,\boldsymbol{\tau}^\top]^\top$；点接触下第 $i$ 块为 $\mathbf{G}_i=[\mathbf{I}_3;\skewm{\mathbf{r}}_i]$（\cref{eq:cm-Gi}），$\mathbf{r}_i$ 是接触点相对质心的位置。\cref{ch:grasp-closure} 已证：\emph{力封闭}$\iff$各摩擦锥经 $\mathbf{G}$ 映出的合力旋量锥充满整个力旋量空间 $\mathbb{R}^6$（平面则 $\mathbb{R}^3$）。本章就在这个“合力旋量锥”上做文章。

% ════════════════════════════════════════════════════════════════
\section{抓取质量度量：从“能否封闭”到“封闭多牢”}
\label{sec:gp-quality}
% ════════════════════════════════════════════════════════════════

\subsection{动机：为何需要一个标量}
\cref{ch:grasp-closure}给出的是\emph{二值}回答：一个抓取要么力封闭、要么不是。但工程中同一物体常有成千上万种力封闭抓取（夹爪可落在物体表面任意一对位置），机器人必须\emph{排序}并选最好的一个。\citeauthor{lynch2017modern}\cite{lynch2017modern}用一句话点破需求：图\,12.16 里两个抓取都形封闭，“哪个更好？”——回答它，必须有一个\textbf{抓取质量度量（grasp metric）}
\begin{equation}
  \mathrm{Qual}(\{\mathbf{G}_i,\,\text{摩擦锥}_i\})\in\mathbb{R},
  \label{eq:gp-qual-map}
\end{equation}
约定 $\mathrm{Qual}<0$ 表示该抓取\emph{不}力封闭，$\mathrm{Qual}>0$ 越大越好。本节给出三个互补的度量：最大内接球、Ferrari--Canny $\varepsilon$、面向任务度量。它们的共同出发点是同一个集合——抓取力旋量集。

\subsection{抓取力旋量集（grasp wrench set）}
\begin{definition}[抓取力旋量集]\label{def:gp-gws}
设第 $i$ 个接触的接触力受限于摩擦锥，且其法向力幅值上界为 $f_{i,\max}>0$（避免压坏物体或受执行器限幅）。则\emph{抓取力旋量集}是所有这些受限接触力经抓取映射合成的物体力旋量之全体：
\begin{equation}
  \mathcal{W}=\Big\{\,\mathbf{w}=\mathbf{G}\mathbf{f}\ \Big|\ \mathbf{f}_i\in\mathrm{FC}_i,\ \|\mathbf{f}_i\|\le f_{i,\max},\ i=1,\dots,N\,\Big\}\subset\mathbb{R}^6,
  \label{eq:gp-gws}
\end{equation}
其中 $\mathrm{FC}_i$ 是第 $i$ 个接触的摩擦锥（\cref{def:friction-cone}）。
\end{definition}

把每个圆摩擦锥用 $m$ 条棱的多面体锥内接近似（\cref{ch:grasp-closure} 的棱锥近似），单接触可施加的力旋量是有限条\emph{基本力旋量} $\{\mathbf{w}_{ij}\}_{j=1}^m$（每条棱经 $\mathbf{G}_i$ 映射、并乘以幅值上界）的凸组合；$N$ 个接触叠加后，$\mathcal{W}$ 是这些基本力旋量的\textbf{凸包}（convex hull）
\begin{equation}
  \mathcal{W}=\conv\big(\{\,\mathbf{w}_{ij}\,\}_{i=1,\dots,N;\ j=1,\dots,m}\big),
  \label{eq:gp-gws-hull}
\end{equation}
是一个凸多胞体（polytope）。这里出现两种约定，务必区分（呼应 \cref{ch:grasp-closure} 与 \cite{ferrari1992}）：
\begin{itemize}
  \item \textbf{$L_\infty$ 约定}：每个接触\emph{独立}限幅 $\|\mathbf{f}_i\|\le f_{i,\max}$。物理上对应“每根手指有独立的力上限”。
  \item \textbf{$L_1$ 约定}：限制\emph{总}接触力 $\sum_i\|\mathbf{f}_i\|\le f_{\max}$。物理上对应“总抓持力/总功率受限”，更贴近单一驱动源经差动机构分配的手。
\end{itemize}
两种约定给出不同形状的 $\mathcal{W}$，因而下文的度量值不可跨约定比较。

\begin{insight}[$\mathbf{G}$ 的几何在度量里复活]
$\mathcal{W}$ 的“胖瘦”直接编码抓取的好坏。力封闭 $\iff$ $\mathbf{0}$ 是 $\mathcal{W}$ 的\emph{内点}（origin in interior）——因为这意味着任意小的外扰力旋量都能被某组接触力顶住。\cref{ch:grasp-closure} 的“合力旋量锥充满 $\mathbb{R}^6$”，在限幅后正是“$\mathcal{W}$ 包含原点的一个邻域”。于是“封闭多牢”自然量化为：原点到 $\mathcal{W}$ \emph{边界}有多远。
\end{insight}

\subsection{度量一：最大内接球（largest-minimum-resisted-wrench）}
\begin{definition}[最大内接球半径度量]\label{def:gp-ball}
$\mathrm{Qual}(\cdot)$ 取为以力旋量空间原点为心、可放入凸多胞体 $\mathcal{W}$ 的\emph{最大球}之半径：
\begin{equation}
  Q_{\text{ball}}=\max\{\,r\ge 0\ :\ \mathcal{B}_r=\{\mathbf{w}:\|\mathbf{w}\|\le r\}\subseteq\mathcal{W}\,\}.
  \label{eq:gp-ball}
\end{equation}
\end{definition}

\begin{proposition}[内接球半径 $=$ 最差方向的可抗力旋量]\label{prop:gp-ball-meaning}
设 $\mathcal{W}$ 由超平面族 $\{\mathbf{w}:\mathbf{n}_k^\top\mathbf{w}\le b_k,\ b_k>0\}$ 的交给出（$\mathbf{n}_k$ 单位外法向，原点在内故 $b_k>0$）。则
\begin{equation}
  Q_{\text{ball}}=\min_k\, b_k,
  \label{eq:gp-ball-faces}
\end{equation}
即\emph{原点到最近一张面的距离}。几何上 $Q_{\text{ball}}$ 等于“在所有单位方向 $\mathbf{u}$ 上，抓取能抵抗的外力旋量幅值的\emph{最小值}”——故称\textbf{最大化最小可抗力旋量}（largest-minimum-resisted-wrench）。
\end{proposition}
\begin{proof}
球 $\mathcal{B}_r\subseteq\mathcal{W}$ 当且仅当对每张面 $\mathbf{n}_k^\top\mathbf{w}\le b_k$ 都有 $\max_{\|\mathbf{w}\|\le r}\mathbf{n}_k^\top\mathbf{w}=r\le b_k$（支撑函数取在 $\mathbf{w}=r\mathbf{n}_k$）。故最大可行 $r=\min_k b_k$。又沿任意单位方向 $\mathbf{u}$，抓取可抗的最大幅值是 $\max\{t:t\mathbf{u}\in\mathcal{W}\}=\min_k b_k/(\mathbf{n}_k^\top\mathbf{u})_+$；其在所有 $\mathbf{u}$ 上的下确界正是 $\min_k b_k$。$\hfill\square$
\end{proof}

\begin{note}[两个必须当心的“单位/原点”陷阱]
\citeauthor{lynch2017modern}\cite{lynch2017modern}明确警告（§12.1.7.3）：（i）\emph{力与力矩量纲不同}，把 $\mathbf{w}=[\mathbf{F};\boldsymbol{\tau}]$ 直接取欧氏范数无物理意义——通常引入物体的\textbf{特征长度} $\ell$（如包围盒尺度），将力矩归一为 $\boldsymbol{\tau}/\ell$，使二者同量纲；（ii）\emph{力矩依赖参考原点}——同一抓取换一个力矩参考点，$\mathcal{W}$ 形状随之改变，度量值也变，故必须固定在物体质心（与 \cref{eq:cm-grasp} 的 $\mathbf{G}$ 参考一致）。不处理这两点，跨物体/跨抓取的数值不可比。
\end{note}

\subsection{度量二：Ferrari--Canny $\varepsilon$ 度量}
最大内接球度量在文献中以 \citeauthor{ferrari1992}\cite{ferrari1992} 的形式最为标准，称 $\varepsilon$ \textbf{度量}（epsilon metric）。其表述等价于 \cref{def:gp-ball}，但强调“最坏情况外扰”的对偶视角，并区分 $L_1/L_\infty$ 约定。

\begin{definition}[Ferrari--Canny $\varepsilon$ 度量]\label{def:gp-epsilon}
\begin{equation}
  \varepsilon=\min_{\|\mathbf{u}\|=1}\ \max\{\,\lambda\ge 0\ :\ \lambda\,\mathbf{u}\in\mathcal{W}\,\}
  \;=\;\min_{\mathbf{w}\in\partial\mathcal{W}}\|\mathbf{w}\|,
  \label{eq:gp-epsilon}
\end{equation}
即\emph{在所有单位外扰方向中，抓取能平衡的最大外扰力旋量幅值之最小者}；等价于原点到 $\mathcal{W}$ 边界 $\partial\mathcal{W}$ 的最近距离。$\varepsilon>0\iff$ 力封闭；$\varepsilon$ 越大抓取越鲁棒。
\end{definition}

\begin{insight}[$\varepsilon$ 是“抗最坏外扰”的保证]
$\varepsilon$ 给的是\emph{最坏情况保证}：无论外扰来自哪个方向，只要其力旋量幅值 $\le\varepsilon$，抓取（在限幅内）就一定能顶住。这正是“最小可抗力旋量最大化”的鲁棒优化解读——它对“我们不知道外扰会从哪来”这一无知，给出可证明的安全裕度。$\varepsilon$ 度量因此成为抓取综合（grasp synthesis）与数据驱动抓取（如 Dex-Net 给抓取打分）最常用的解析标签。
\end{insight}

\begin{example}[二维三指抓取的 $\varepsilon$]\label{ex:gp-epsilon-2d}
平面物体（力旋量空间 $\mathbb{R}^3$，$\mathbf{w}=[m_z,f_x,f_y]^\top$）受三个带摩擦点接触。每个摩擦锥贡献两条棱、各映成一条基本力旋量，共 $6$ 条 $\mathbf{w}_{ij}\in\mathbb{R}^3$。取 $f_{i,\max}=1$，$\mathcal{W}=\conv\{\mathbf{w}_{ij}\}$ 是 $\mathbb{R}^3$ 中一个凸多胞体。若原点在其内部，逐面算 $b_k=\mathbf{n}_k^\top$（面到原点的有向距离），则 $\varepsilon=\min_k b_k>0$。直觉：三指越“对称包围”物体、摩擦系数越大，$\mathcal{W}$ 越接近以原点为心的球，$\varepsilon$ 越大；三指挤在物体一侧则某个 $b_k\to 0$，$\varepsilon\to 0$（接近失封闭）。
\end{example}

\subsection{度量三：面向任务的力旋量空间度量}
最大内接球/$\varepsilon$ 度量隐含“外扰各向同性”（用球衡量）。但真实任务的外扰常\emph{高度各向异性}：拧瓶盖主要需绕轴力矩，端水平托盘主要需竖直力。此时应把“球”换成\textbf{任务力旋量空间}（task wrench space, TWS）——任务可能施加的外扰力旋量集合 $\mathcal{T}\subset\mathbb{R}^6$（如由任务采样或专家示范统计得到）。

\begin{definition}[面向任务的抓取度量]\label{def:gp-task}
给定任务力旋量集 $\mathcal{T}$（含原点、凸、对称归一），定义
\begin{equation}
  Q_{\text{task}}=\max\{\,s\ge 0\ :\ s\,\mathcal{T}\subseteq\mathcal{W}\,\},
  \label{eq:gp-task}
\end{equation}
即把任务力旋量集放大到刚好仍被抓取力旋量集容纳的\emph{最大缩放因子}。$Q_{\text{task}}\ge 1\iff$ 抓取能完成该任务的全部外扰需求；$Q_{\text{task}}$ 越大越有裕度。
\end{definition}

\begin{note}[三个度量的统一与退化]
三者是同一思想（“$\mathcal{W}$ 能容纳多大的需求集”）的不同需求集选择：
\[
\underbrace{\mathcal{B}_1\ (\text{单位球})}_{Q_{\text{ball}}=\varepsilon}\quad\subset\quad\underbrace{\text{任务集}\ \mathcal{T}}_{Q_{\text{task}}}.
\]
取 $\mathcal{T}=\mathcal{B}_1$（各向同性球）则 $Q_{\text{task}}=\varepsilon$，故 $\varepsilon$ 是面向任务度量的\emph{无任务先验}特例。工程上：无任务信息时用 $\varepsilon$（保守、通用）；有任务先验时用 $Q_{\text{task}}$（更省力、更贴合）。
\end{note}

% ════════════════════════════════════════════════════════════════
\section{抓取规划：对踵抓取、接触数界与构造}
\label{sec:gp-planning}
% ════════════════════════════════════════════════════════════════

上一节回答“给定一个抓取，多好？”；本节回答“\emph{怎样选}接触点，使抓取力封闭（并尽量好）？”。这就是\textbf{抓取规划}（grasp planning）。我们从两个理论界（需要几个接触）入手，再给二指情形的完整充要判据与构造算法。

\subsection{接触数界：Caratheodory 与 Steinitz}
\cref{ch:grasp-closure} 已证：无摩擦点接触下，力封闭 $\iff$ 各接触法向力旋量\emph{正张成}（positively span）整个力旋量空间 $\mathbb{R}^p$（$p=3$ 平面、$p=6$ 空间）。两个凸分析的经典定理把“正张成 $\mathbb{R}^p$ 需几个向量”框死。

\begin{theorem}[Caratheodory：正张成下界]\label{thm:gp-caratheodory}
若向量集 $X=\{\mathbf{v}_1,\dots,\mathbf{v}_k\}$ 正张成 $\mathbb{R}^p$，则 $k\ge p+1$。
\end{theorem}

\begin{theorem}[Steinitz：正张成上界]\label{thm:gp-steinitz}
若 $S\subset\mathbb{R}^p$ 且 $\mathbf{q}\in\interior(\conv S)$，则存在 $X=\{\mathbf{v}_1,\dots,\mathbf{v}_k\}\subseteq S$ 使 $\mathbf{q}\in\interior(\conv X)$ 且 $k\le 2p$。
\end{theorem}

\begin{proposition}[无摩擦点接触的力封闭接触数界]\label{prop:gp-contact-bounds}
对\emph{非例外}（non-exceptional）物体，用无摩擦点接触实现力封闭所需接触数 $k$ 满足
\begin{equation}
  p+1\ \le\ k\ \le\ 2p,\qquad
  \begin{cases}
    4\le k\le 6, & p=3\ (\text{平面}),\\
    7\le k\le 12, & p=6\ (\text{空间}).
  \end{cases}
  \label{eq:gp-bounds}
\end{equation}
\end{proposition}
\begin{proof}[证明思路（\cite{murray1994mathematical}）]
下界：力封闭要求 $\mathbf{0}$ 被接触法向力旋量\emph{正}张成（即原点在凸包内部），由 \cref{thm:gp-caratheodory} 需 $k\ge p+1$。上界：将 $S$ 取为物体边界上所有可用法向力旋量之集、$\mathbf{q}=\mathbf{0}$；若该抓取力封闭则 $\mathbf{0}\in\interior(\conv S)$，由 \cref{thm:gp-steinitz} 可从中选出 $\le 2p$ 个仍使 $\mathbf{0}$ 在其凸包内部，即存在 $\le 2p$ 接触的力封闭抓取。$\hfill\square$
\end{proof}

\begin{note}[例外曲面：旋转面永远抓不住]
\cref{eq:gp-bounds} 的前提是物体\emph{非例外}。\textbf{例外曲面}（exceptional surface）指其上所有点的法向力旋量之凸包\emph{不}含原点邻域的曲面——典型是\emph{旋转面}：球、椭球、圆柱、圆盘。直觉：球面任一点的法线都过球心，绕球心的任意力矩\emph{无人能抗}，故纯无摩擦点接触永远无法力封闭一个球（无论几个接触）。这与 \cref{ch:grasp-closure} 形封闭里“可绕中心自旋的圆盘”是同一现象的力学版本。例外物体必须靠\emph{摩擦}（指尖切向力提供绕轴力矩）才能抓住。
\end{note}

\begin{insight}[摩擦把“$7$ 个”降到“$3$ 个”]
\cref{ch:grasp-closure} 说空间形封闭至少 $7$ 个点接触，本节说无摩擦力封闭也要 $7\sim 12$ 个。但真实手指\emph{有摩擦}：每个带摩擦点接触不再只出一条法向力旋量，而是出一整个\emph{摩擦锥}（$3$ 维力的连续族）。摩擦锥极大丰富了可施加力旋量集，使空间力封闭仅需 $3$ 个带摩擦接触、平面仅需 $2$ 个（下一小节）。代价是判据从“正张成 $\iff$ 凸包含原点”变为“摩擦锥的合成锥含原点邻域”（\cref{ch:grasp-closure} 的力封闭 LP/凸包判据）。\textbf{结论}：手指数的锐减，本质是\emph{摩擦}买来的。
\end{insight}

\subsection{二指对踵抓取定理}
最常用的机器人末端是平行夹爪——\emph{两}个接触。两指能否力封闭，有一条极其简洁的几何充要条件。

\begin{theorem}[平面对踵抓取（antipodal grasp）\cite{nguyen1988}]\label{thm:gp-antipodal}
平面物体被两个带摩擦点接触抓取，\emph{力封闭}当且仅当连接两接触点的直线\emph{同时落在两个接触的摩擦锥内}。满足此条件的抓取称\textbf{对踵抓取}。
\end{theorem}

\begin{proposition}[空间软指对踵推广\cite{nguyen1988}]\label{prop:gp-antipodal-spatial}
空间物体被两个\emph{软指接触}（soft-finger，传 $3$ 维力 $+$ 绕法向摩擦力矩）抓取，力封闭当且仅当连接两接触点的直线落在两个摩擦锥内。
\end{proposition}

\begin{derivation}[对踵条件的几何直觉]
设两接触点为 $\mathbf{p}_1,\mathbf{p}_2$，内法向 $\mathbf{n}_1,\mathbf{n}_2$。“连线在两锥内”意味着：手指 $1$ 可沿\emph{指向 $\mathbf{p}_2$ 的方向}施力（该方向在锥 $1$ 内），手指 $2$ 可沿\emph{指向 $\mathbf{p}_1$ 的方向}施力。于是两指可施加一对\emph{共线、反向、过两点连线}的力——它们合力为零、合力矩为零，构成纯\textbf{内力}（挤压），把物体“夹死”在两指之间且不产生净力旋量。任何外扰力旋量都能用“调节内力大小 $+$ 在锥内偏转接触力”来平衡：因为每个锥都\emph{张开成一个 $2$ 维（平面）扇区}，两扇区合起来能生成 $\mathbb{R}^3$ 中任意力旋量方向。反之若连线\emph{出}了某个锥，则存在一个绕连线的外力矩方向，两接触力（都被困在各自锥内、无法产生足够反向力矩）无法抵抗——失封闭。这正是 \cref{ch:grasp-closure} 中“两锥能互相‘看见’对方底面”的几何刻画。$\hfill\square$
\end{derivation}

\begin{note}[“对踵”一词的来历与边界]
antipodal $=$“对跖/正对”，指两接触点处的法向近乎\emph{正对}（连线落在锥内即要求法向偏离连线不超过摩擦角 $\alpha=\arctan\mu$）。摩擦越大、锥越宽，可接受的“非正对”越多。务必注意：对踵是\emph{二指}专属充要；三指及以上不能化简为如此简洁的几何条件，须回到 \cref{ch:grasp-closure} 的合成锥/LP 判据。
\end{note}

\subsection{构造力封闭抓取的算法思想}
\cref{thm:gp-antipodal} 不只是判据，还直接给出\emph{平面多边形}上枚举全部力封闭抓取的算法\cite{murray1994mathematical}。

\begin{definition}[对踵条件的线性不等式形式]\label{def:gp-antipodal-ineq}
设 $\mathbf{n}_i$ 为接触点 $\mathbf{p}_i$ 处内法向，$\mathbf{n}^\perp$ 为 $\mathbf{n}$ 逆时针转 $90^\circ$ 的单位向量，$\otimes$ 为平面叉积 $\mathbf{a}\otimes\mathbf{b}=\det[\mathbf{a}\ \mathbf{b}]$。则“连线 $\mathbf{p}_2-\mathbf{p}_1$ 落在两摩擦锥内”等价于四条线性（对固定边对）不等式：
\begin{equation}
\begin{aligned}
  (\mathbf{n}_1-\mu\mathbf{n}_1^\perp)\otimes(\mathbf{p}_2-\mathbf{p}_1)&\ge 0, &
  (\mathbf{n}_1+\mu\mathbf{n}_1^\perp)\otimes(\mathbf{p}_2-\mathbf{p}_1)&\le 0,\\
  (\mathbf{n}_2-\mu\mathbf{n}_2^\perp)\otimes(\mathbf{p}_1-\mathbf{p}_2)&\ge 0, &
  (\mathbf{n}_2+\mu\mathbf{n}_2^\perp)\otimes(\mathbf{p}_1-\mathbf{p}_2)&\le 0.
\end{aligned}
  \label{eq:gp-antipodal-ineq}
\end{equation}
\end{definition}

\begin{algo}[平面多边形力封闭抓取的枚举与“最大独立接触区”\cite{murray1994mathematical}]\label{algo:gp-construct}
\hfill
\begin{enumerate}
  \item \textbf{枚举边对}：对多边形每一对边 $(e_a,e_b)$，把接触点 $\mathbf{p}_1\in e_a,\mathbf{p}_2\in e_b$ 参数化为边上弧长。
  \item \textbf{解不等式区}：在 $(\mathbf{p}_1,\mathbf{p}_2)$ 参数平面上，把 \cref{eq:gp-antipodal-ineq} 四条不等式各取等号画成直线，四条不等式的公共可行域即该边对上\emph{全部}力封闭抓取的集合（一个凸多边形区域）。
  \item \textbf{并集}：对所有边对取并，得物体全部二指力封闭抓取。
  \item \textbf{择优}：在可行域内选最优抓取。两种常用准则——
  （a）最大化某抓取质量度量（\cref{sec:gp-quality} 的 $\varepsilon$）；
  （b）选\emph{离可行域边界最远}的点，等价于在每条边上找一对\textbf{最大独立接触区}（maximal independent regions）：只要两接触各自落在其区内，抓取必力封闭——对接触定位误差最鲁棒（可行域内最大内接方块的两条边投影）。
\end{enumerate}
\end{algo}

\begin{insight}[“判据—不等式—枚举”三步走是抓取规划的通法]
\cref{algo:gp-construct} 的范式可推广：把\emph{封闭判据}写成接触参数的\emph{（线性或非线性）不等式}，再在接触参数空间\emph{枚举/采样/优化}求可行域。曲面物体只需把 \cref{eq:gp-antipodal-ineq} 换成接触位置的非线性函数（用曲面参数表达法向），可行域求解从“画直线”变为“求零集”。这与数据驱动抓取（采样大量候选抓取、用 $\varepsilon$ 或学习模型打分筛选）是同一思想的连续优化版与采样版。
\end{insight}

% ════════════════════════════════════════════════════════════════
\section{非抓持与准静态操作：推动、倾倒、动态抓取}
\label{sec:gp-nonprehensile}
% ════════════════════════════════════════════════════════════════

至此都在谈\emph{抓持}（prehensile，把物体锁在手内）。但操作远不止抓取——\citeauthor{lynch2017modern}\cite{lynch2017modern}列举：托盘端杯、绕脚翻冰箱、推沙发、抛接球、振动送料……这些\textbf{非抓持操作}（nonprehensile manipulation）让机器人能动“抓不住/太重/太大”的物体。本节用统一框架——接触运动学（\cref{ch:grasp-closure}）$+$ 库仑摩擦 $+$ 刚体动力学——分析其中三类：推动、倾倒/翻转、动态抓取。

\subsection{统一框架：接触模式枚举与拟静态约化}
单刚体在多个接触下的运动由下式（物体质心系，记号承 \cref{ch:arm-dyn} 的刚体动力学）确定\cite{lynch2017modern}：
\begin{equation}
  \mathbf{w}_{\text{ext}}+\sum_i k_i\,\mathbf{w}_i=\mathbf{G}_{\text{inertia}}\dot{\boldsymbol{\mathcal{V}}}-[\mathrm{ad}_{\boldsymbol{\mathcal{V}}}]^\top\mathbf{G}_{\text{inertia}}\boldsymbol{\mathcal{V}},
  \quad k_i\ge 0,\ \mathbf{w}_i\in\text{接触力旋量锥}_i,
  \label{eq:gp-rb-dyn}
\end{equation}
其中 $\boldsymbol{\mathcal{V}}$ 是物体 twist、$\mathbf{G}_{\text{inertia}}$ 是空间惯量矩阵、$\mathbf{w}_{\text{ext}}$ 是重力等外力旋量、$k_i\mathbf{w}_i$ 是第 $i$ 接触的接触力旋量。求解方法是\textbf{接触模式枚举}：

\begin{algo}[刚体接触运动的接触模式分析\cite{lynch2017modern}]\label{algo:gp-contact-mode}
\hfill
\begin{enumerate}
  \item \textbf{列举接触模式}：每个接触贴标签 $R$（粘着/滚动 rolling-sticking）、$S$（滑动 sliding）、$B$（分离 breaking），物体的接触模式是各接触标签的串接。
  \item \textbf{逐模式判一致性}：对每种模式，问是否存在接触力旋量 $\sum_i k_i\mathbf{w}_i$（满足库仑律与该模式标签）与加速度 $\dot{\boldsymbol{\mathcal{V}}}$（满足该模式的运动学约束），使 \cref{eq:gp-rb-dyn} 成立。若有，该“模式 $+$ 接触力 $+$ 加速度”是一个\emph{一致解}。
\end{enumerate}
\end{algo}

\begin{note}[刚体 $+$ 库仑摩擦的代价：歧义与不一致]
\cref{algo:gp-contact-mode} 是“逐情形分析”而非解一组方程，且可能找到\emph{多个}一致解（\textbf{歧义问题}，ambiguous）或\emph{零个}一致解（\textbf{不一致问题}，inconsistent）。这令人不安——现实力学问题理应唯一可解！但这正是“理想刚体 $+$ 库仑摩擦”假设的代价（与 \cref{ch:grasp-closure} 中刚体接触的固有病态同源）。所幸多数实际问题该法仍给唯一模式与运动。
\end{note}

\begin{definition}[拟静态（quasistatic）操作]\label{def:gp-quasistatic}
若运动足够慢，惯性力可略，\cref{eq:gp-rb-dyn} 退化为\textbf{力旋量平衡}：
\begin{equation}
  \mathbf{w}_{\text{ext}}+\sum_i k_i\,\mathbf{w}_i=\mathbf{0},\qquad k_i\ge 0,\ \mathbf{w}_i\in\text{接触力旋量锥}_i.
  \label{eq:gp-quasistatic}
\end{equation}
即任何时刻接触力旋量与外力旋量精确平衡。这是分析推动、装配、堆叠稳定性的主力工具。
\end{definition}

\subsection{推动的运动学与摩擦极限面}
\label{sec:gp-pushing}
推动是非抓持操作的原型：一根手指（或一条边）推地上的物体，物体在\emph{支撑面摩擦}与\emph{推力}的合力下平动 $+$ 转动。难点在于：地面对物体底面的分布式摩擦，如何合成为一个作用于物体的力旋量？答案是\textbf{摩擦极限面}（limit surface）。

\begin{definition}[摩擦极限面\cite{goyal1991planar}]\label{def:gp-limit-surface}
物体底面与支撑面间的分布式库仑摩擦，对物体质心产生一个平面力旋量 $\mathbf{w}=[f_x,f_y,m_z]^\top$。当底面相对支撑面以某\emph{瞬时运动旋量}（平面 twist，绕某转动中心 CoR）滑动时，所有可能的摩擦力旋量构成 $\mathbf{w}$ 空间中一张闭凸曲面，称\textbf{极限面} $\mathcal{L}$。其关键性质（最大耗散原理的推论）：
\begin{equation}
  \boxed{\;\text{物体的滑动 twist}\ \boldsymbol{\mathcal{V}}\ \text{方向}\ \perp\ \text{极限面}\ \mathcal{L}\ \text{在对应摩擦力旋量点处的切平面}\;}
  \label{eq:gp-limit-normality}
\end{equation}
即\emph{滑动 twist 沿极限面的外法向}（normality / maximum-dissipation 法则）。
\end{definition}

\begin{insight}[极限面把“分布摩擦”降成“一张曲面查表”]
极限面的威力在于：它把支撑面上\emph{无穷多}个接触点的分布摩擦，浓缩为质心力旋量空间中\emph{一张曲面}。给定外推力旋量 $\mathbf{w}$（落在 $\mathcal{L}$ 上），\cref{eq:gp-limit-normality} 的法向法则立刻读出物体的滑动 twist 方向（绕哪点、怎么转）——无须对底面积分。这正是推动规划（push planning）与触觉重定位（tactile re-orientation）的核心工具。实务中常把 $\mathcal{L}$ 近似为一个\textbf{椭球}（ellipsoidal limit surface），则法向法则化为简单的线性映射 $\boldsymbol{\mathcal{V}}\propto\mathbf{A}\mathbf{w}$。
\end{insight}

\begin{example}[拟静态推动的转动中心\cite{lynch2017modern,goyal1991planar}]\label{ex:gp-push}
一根手指推一个平放盒子，盒子慢速运动（拟静态）。两个力旋量平衡（\cref{eq:gp-quasistatic}）：推力旋量 $+$ 地面极限面摩擦力旋量 $=\mathbf{0}$。于是地面须提供一个与推力旋量反向的摩擦力旋量；该力旋量在极限面 $\mathcal{L}$ 上的点，其外法向即盒子的滑动 twist——决定盒子绕\emph{哪个转动中心}转、向左偏还是向右偏。\textbf{关键洞见}：盒子转向由\emph{支撑摩擦分布的几何}（极限面形状）与\emph{推力作用线}共同决定，而非仅由质量分布。这解释了“推一个物体它常常一边走一边转”的日常现象。
\end{example}

\subsection{倾倒、翻转与拟静态稳定性}
拟静态平衡同样判定\emph{多体堆叠}能否站立、以及如何\emph{翻转}（pivoting/tumbling）物体。

\begin{example}[拱的稳定性\cite{lynch2017modern}]\label{ex:gp-arch}
一座由数块石头叠成的拱，在重力下是否稳定？对每块石 $j$ 写拟静态力旋量平衡（\cref{eq:gp-quasistatic}）：$\sum_{i}k_i\mathbf{w}_i+\mathbf{w}_{\text{ext},j}=\mathbf{0}$，且相邻石块间作用力等大反向。这是一组以 $k_i\ge 0$ 为变量的\textbf{线性约束满足问题}（可由线性规划求解）：若存在一组非负 $k_i$ 使所有石块力旋量平衡（且各接触力在摩擦锥内），拱稳定；否则坍塌。\emph{倾倒/翻转}则是“让平衡\emph{临界破坏}”：缓慢移动接触/施加外力，使物体绕某接触棱的力矩平衡恰被打破，物体绕该棱翻转一个面。
\end{example}

\subsection{动态抓取：用惯性力代替封闭}
当不满足任何静态封闭、却仍想搬运物体时，可借\emph{惯性力}。

\begin{example}[动态抓取（dynamic grasp）：加速度顶住物体\cite{lynch2017modern}]\label{ex:gp-dynamic}
平面方块被两指托住，一指 $\mu=1$、另一指无摩擦。\emph{静止}时（\cref{eq:gp-quasistatic}），两指能施加的合成力旋量锥\emph{不含}平衡重力所需的力旋量——方块会相对手指滑动，纯静态托不住。但若两指\emph{一起向左加速} $2g$，则在“方块随手指同步运动（接触模式 $RR$）”假设下，方块也须以 $2g$ 向左加速；产生此加速度所需的力旋量 $+$ 平衡重力的力旋量，其矢量和\emph{落入}了合成力旋量锥（\cref{eq:gp-rb-dyn} 含惯性项）。于是 $RR$ 成为一致解：\textbf{惯性力把方块压向手指}，手指无须封闭即可“握住”运动中的方块。代数上对加速度 $a_x$ 解 $k_1,k_2,k_3\ge 0$，得可行区间（如 $-5g\le a_x\le -g$）。\textbf{要点}：动态抓取靠\emph{运动}维持接触，故须确保该运动下\emph{无其他接触模式}（滑动/分离）也一致，否则物体可能甩脱。
\end{example}

\begin{example}[米尺翻手（meter-stick）：拟静态滑移的自校正\cite{lynch2017modern}]\label{ex:gp-meterstick}
把米尺横搁两食指上，缓慢相向移动两指。直觉以为某指越过质心时尺会掉——实则不会。拟静态分析（\cref{eq:gp-quasistatic}）：三种接触模式 $RS_r$（左指粘、右指滑）、$S_lR$（左指滑、右指粘）、$S_lS_r$（两指都滑）中，\emph{只有 $S_lR$} 能提供平衡重力的力旋量——即\emph{承重更多的那根指（离质心近者）保持不滑，另一指在尺下滑动}。于是质心始终被“拽”向两指中点，直到两指相遇尺也不掉。这是非抓持操作里“摩擦自动分配、系统自校正”的优美例子；它依赖拟静态假设——若快速移指，惯性需求超出粘着摩擦上限，尺就会掉。
\end{example}

\begin{example}[峰嵌（wedging）与卡阻（jamming）：装配的力封闭副产物\cite{lynch2017modern}]\label{ex:gp-wedge}
力控插销与孔两点接触。若控制器施加的力旋量 $\mathbf{w}_1$ 能被两接触摩擦锥\emph{平衡}，销\textbf{卡阻}（jam）停住；若施加 $\mathbf{w}_2$ 不能被平衡，则插入继续。更甚者，若两接触摩擦锥大到能“互相看见对方底面”（即 \cref{thm:gp-antipodal} 的对踵条件在孔内成立），销处于\emph{力封闭}状态、可抗任意力旋量——称\textbf{峰嵌}（wedged），无论怎么推都卡死。这把本章的力封闭判据反向用到了\emph{装配避卡}：插孔策略要\emph{避开}力封闭/卡阻的力旋量方向。
\end{example}

% ════════════════════════════════════════════════════════════════
\section{滚动接触运动学：Montana 接触方程}
\label{sec:gp-rolling}
% ════════════════════════════════════════════════════════════════

真实的多指操作里，手指是\emph{曲面}，在物体曲面上\emph{滚动}（而非固定点接触）。要做手内操作（in-hand manipulation）——靠滚动重定位物体——必须知道：当两曲面相对运动时，\emph{接触点在各自曲面上如何爬行}。这由 \citeauthor{montana1988}\cite{montana1988} 的接触方程精确刻画，是 MLS Ch5 的理论高峰\cite{murray1994mathematical}。本节从微分几何（第一/二基本形式）出发，推到 Montana 方程，并衔接 \cref{ch:lie} 的 $\SE(3)$ 旋量。

\subsection{曲面的局部几何：第一/二基本形式与曲率张量}
设曲面 $S$ 由局部正交坐标卡 $\mathbf{c}:U\subset\mathbb{R}^2\to\mathbb{R}^3$，$(u,v)\mapsto\mathbf{c}(u,v)$ 描述（在物体系下）。切平面由 $\mathbf{c}_u=\partial\mathbf{c}/\partial u$、$\mathbf{c}_v=\partial\mathbf{c}/\partial v$ 张成。

\begin{definition}[第一基本形式与度量张量\cite{murray1994mathematical}]\label{def:gp-first-form}
\emph{第一基本形式}给出切向量内积的矩阵
\begin{equation}
  \mathbf{I}_p=\begin{bmatrix}\mathbf{c}_u^\top\mathbf{c}_u & \mathbf{c}_u^\top\mathbf{c}_v\\ \mathbf{c}_v^\top\mathbf{c}_u & \mathbf{c}_v^\top\mathbf{c}_v\end{bmatrix},
  \label{eq:gp-first-form}
\end{equation}
对称正定；正交卡下为对角阵。\emph{度量张量} $\mathbf{M}_p$ 是其“平方根”：满足 $\mathbf{I}_p=\mathbf{M}_p^\top\mathbf{M}_p$；正交卡下 $\mathbf{M}_p=\diagop(\|\mathbf{c}_u\|,\|\mathbf{c}_v\|)$，用于把坐标速度归一化为切平面上的物理速度。
\end{definition}

\begin{definition}[Gauss 映射、第二基本形式与曲率张量\cite{murray1994mathematical}]\label{def:gp-second-form}
\emph{Gauss 映射} $\mathbf{n}=N(u,v)=\dfrac{\mathbf{c}_u\times\mathbf{c}_v}{\|\mathbf{c}_u\times\mathbf{c}_v\|}$ 给出单位外法向。\emph{第二基本形式}度量法向沿曲面的变化（即曲率）：
\begin{equation}
  \mathbf{II}_p=\begin{bmatrix}\mathbf{c}_u^\top\mathbf{n}_u & \mathbf{c}_u^\top\mathbf{n}_v\\ \mathbf{c}_v^\top\mathbf{n}_u & \mathbf{c}_v^\top\mathbf{n}_v\end{bmatrix},\quad \mathbf{n}_u=\partial\mathbf{n}/\partial u,\ \mathbf{n}_v=\partial\mathbf{n}/\partial v.
  \label{eq:gp-second-form}
\end{equation}
归一化后得\textbf{曲率张量}（curvature form）
\begin{equation}
  \mathbf{K}_p=\mathbf{M}_p^{-\top}\,\mathbf{II}_p\,\mathbf{M}_p^{-1}\ \in\mathbb{R}^{2\times 2},
  \label{eq:gp-curvature}
\end{equation}
它是“在归一化 Gauss 标架下、单位法向随单位切向位移的变化率”，对称（特征值为两个主曲率）。平面 $\mathbf{K}_p=\mathbf{0}$。
\end{definition}

\begin{definition}[归一化 Gauss 标架与挠率形式\cite{murray1994mathematical}]\label{def:gp-gauss-frame}
正交卡下定义\textbf{归一化 Gauss 标架}（右手系，$z$ 轴取外法向）
\begin{equation}
  [\,\mathbf{x}\ \ \mathbf{y}\ \ \mathbf{z}\,]=\Big[\ \tfrac{\mathbf{c}_u}{\|\mathbf{c}_u\|}\ \ \tfrac{\mathbf{c}_v}{\|\mathbf{c}_v\|}\ \ \mathbf{n}\ \Big].
  \label{eq:gp-gauss-frame}
\end{equation}
\emph{挠率形式}（torsion form，行向量）刻画 Gauss 标架沿曲面的“扭转”：
\begin{equation}
  \mathbf{T}_p=\mathbf{y}^\top\Big[\ \tfrac{\partial\mathbf{x}}{\partial u}\big/\|\mathbf{c}_u\|\ \ \ \tfrac{\partial\mathbf{x}}{\partial v}\big/\|\mathbf{c}_v\|\ \Big]\ \in\mathbb{R}^{1\times 2}.
  \label{eq:gp-torsion}
\end{equation}
三元组 $(\mathbf{M}_p,\mathbf{K}_p,\mathbf{T}_p)$ 合称曲面的\textbf{几何参数}，是接触方程的全部输入。
\end{definition}

\subsection{单接触标架的诱导速度}
设接触点沿曲面画出曲线 $\mathbf{p}(t)\in S$，接触标架 $C$ 时刻与该点的 Gauss 标架重合，局部坐标 $\boldsymbol{\alpha}(t)=\mathbf{c}^{-1}(\mathbf{p}(t))$。

\begin{lemma}[接触标架的诱导（体）速度\cite{montana1988,murray1994mathematical}]\label{lem:gp-montana-induced}
接触标架 $C$ 相对物体参考标架 $O$ 的体速度旋量（$\SE(3)$ 体 twist，平移在前、记号承 \cref{ch:lie}）由几何参数与坐标速度 $\dot{\boldsymbol{\alpha}}$ 给出：
\begin{equation}
  \mathbf{v}_{oc}=\begin{bmatrix}\mathbf{M}_p\dot{\boldsymbol{\alpha}}\\ 0\end{bmatrix},\qquad
  \boldsymbol{\omega}_{oc}=\begin{bmatrix}-\,(\mathbf{K}_p\mathbf{M}_p\dot{\boldsymbol{\alpha}})_y\\[1pt] (\mathbf{K}_p\mathbf{M}_p\dot{\boldsymbol{\alpha}})_x\\[1pt] \mathbf{T}_p\mathbf{M}_p\dot{\boldsymbol{\alpha}}\end{bmatrix}.
  \label{eq:gp-montana-induced}
\end{equation}
即：平移速度 $=$ 度量张量把坐标速度还原为切平面物理速度；角速度的切向两分量由\emph{曲率张量}决定（标架随曲面弯曲而转），法向分量由\emph{挠率}决定（标架随曲面扭转而绕法向转）。
\end{lemma}

\begin{insight}[几何参数 $=$ 标架运动的“传动比”]
\cref{eq:gp-montana-induced} 揭示 $(\mathbf{M},\mathbf{K},\mathbf{T})$ 的物理角色：当接触点在曲面上以坐标速度 $\dot{\boldsymbol{\alpha}}$ 爬行，$\mathbf{M}$ 决定它\emph{走多快}（平移），$\mathbf{K}$ 决定接触标架\emph{怎么俯仰偏转}（因曲面弯曲），$\mathbf{T}$ 决定它\emph{怎么绕法向自转}（因曲面扭转）。曲面越平（$\mathbf{K},\mathbf{T}\to 0$），接触标架越只平移不转——回到平面情形。
\end{insight}

\subsection{两曲面接触：Montana 接触方程}
现让物体曲面 $S_o$ 与手指曲面 $S_f$ 在一点接触。接触坐标取 $\boldsymbol{\eta}=(\boldsymbol{\alpha}_f,\boldsymbol{\alpha}_o,\psi)$：两曲面各自的局部坐标 $\boldsymbol{\alpha}_f,\boldsymbol{\alpha}_o\in\mathbb{R}^2$，加\textbf{接触角} $\psi$（两接触标架切向轴的夹角）。两曲面相对运动用\emph{在接触点测得}的相对体速度旋量描述：
\begin{equation}
  \boldsymbol{\mathcal{V}}_{\text{rel}}=[\,v_x,v_y,v_z,\ \omega_x,\omega_y,\omega_z\,]^\top,
  \label{eq:gp-rel-twist}
\end{equation}
其中 $(\omega_x,\omega_y)$ 是切平面内的\emph{滚动角速度}、$\omega_z$ 是绕法向的\emph{自旋}、$(v_x,v_y)$ 是切向\emph{滑移速度}、$v_z$ 是法向速度（保持接触则 $v_z=0$）。

\begin{theorem}[Montana 接触方程\cite{montana1988,murray1994mathematical}]\label{thm:gp-montana}
保持接触（$v_z=0$）时，接触坐标 $\boldsymbol{\eta}=(\boldsymbol{\alpha}_f,\boldsymbol{\alpha}_o,\psi)$ 随相对运动旋量按下式演化：
\begin{align}
  \dot{\boldsymbol{\alpha}}_f&=\mathbf{M}_f^{-1}\,(\mathbf{K}_f+\widetilde{\mathbf{K}}_o)^{-1}
    \left(\begin{bmatrix}-\,\omega_y\\[1pt] \omega_x\end{bmatrix}
    -\widetilde{\mathbf{K}}_o\begin{bmatrix}v_x\\ v_y\end{bmatrix}\right),
  \label{eq:gp-montana-f}\\[2pt]
  \dot{\boldsymbol{\alpha}}_o&=\mathbf{M}_o^{-1}\mathbf{R}_\psi\,(\mathbf{K}_f+\widetilde{\mathbf{K}}_o)^{-1}
    \left(\begin{bmatrix}-\,\omega_y\\[1pt] \omega_x\end{bmatrix}
    +\mathbf{K}_f\begin{bmatrix}v_x\\ v_y\end{bmatrix}\right),
  \label{eq:gp-montana-o}\\[2pt]
  \dot{\psi}&=\omega_z+\mathbf{T}_f\mathbf{M}_f\dot{\boldsymbol{\alpha}}_f+\mathbf{T}_o\mathbf{M}_o\dot{\boldsymbol{\alpha}}_o,
  \label{eq:gp-montana-psi}\\[2pt]
  0&=v_z,
  \label{eq:gp-montana-vz}
\end{align}
其中 $\mathbf{R}_\psi=\begin{bmatrix}\cos\psi & -\sin\psi\\ -\sin\psi & -\cos\psi\end{bmatrix}$ 是手指接触标架 $C_f$ 的 $x,y$ 轴相对物体接触标架 $C_o$ 的 $x,y$ 轴的取向，$\widetilde{\mathbf{K}}_o=\mathbf{R}_\psi\mathbf{K}_o\mathbf{R}_\psi$ 是\emph{物体}曲率张量转到手指接触标架切向轴后的表示，$\mathbf{K}_f+\widetilde{\mathbf{K}}_o$ 称\textbf{相对曲率形式}（relative curvature form）。
\end{theorem}

\begin{note}[相对曲率必须非奇异：手内操作的可控前提]
\cref{thm:gp-montana} 只在相对曲率形式\emph{可逆}时成立。奇异的典型：一凸一凹两曲面、曲率半径恰好相等——此时物体微小运动会引起接触点\emph{剧烈}跳变（连续性丧失）。故手内滚动操作须假定运动发生在相对曲率可逆的开集上。这是“能否用滚动稳定地重定位物体”的可控性前提，与 \cref{ch:arm-dualarm-planning} 约束流形的非奇异性要求精神一致。
\end{note}

\begin{derivation}[纯滚动的简化]
\textbf{纯滚动}（rolling without slipping）即无切向滑移、无自旋：$v_x=v_y=0$、$\omega_z=0$。代入 \cref{eq:gp-montana-f,eq:gp-montana-o,eq:gp-montana-psi}，含 $(v_x,v_y)$ 与 $\omega_z$ 的项全消，得
\begin{align}
  \dot{\boldsymbol{\alpha}}_f&=\mathbf{M}_f^{-1}(\mathbf{K}_f+\widetilde{\mathbf{K}}_o)^{-1}\begin{bmatrix}-\omega_y\\ \omega_x\end{bmatrix},
  &\dot{\boldsymbol{\alpha}}_o&=\mathbf{M}_o^{-1}\mathbf{R}_\psi(\mathbf{K}_f+\widetilde{\mathbf{K}}_o)^{-1}\begin{bmatrix}-\omega_y\\ \omega_x\end{bmatrix},
  \label{eq:gp-montana-roll}\\
  \dot{\psi}&=\mathbf{T}_f\mathbf{M}_f\dot{\boldsymbol{\alpha}}_f+\mathbf{T}_o\mathbf{M}_o\dot{\boldsymbol{\alpha}}_o. \notag
\end{align}
此时接触坐标只由\emph{滚动角速度} $(\omega_x,\omega_y)$ 驱动——纯滚动下两接触轨迹等长（“滚过的弧相等”），故 $\dot{\boldsymbol{\alpha}}_f,\dot{\boldsymbol{\alpha}}_o$ 都正比于滚动角速度。$\hfill\square$
\end{derivation}

\begin{example}[球在平面上纯滚动\cite{murray1994mathematical}]\label{ex:gp-sphere-plane}
半径 $\rho$ 的球形指尖在平面物体上纯滚（图省）。平面：$\mathbf{K}_o=\mathbf{0}$、$\mathbf{T}_o=\mathbf{0}$、$\mathbf{M}_o=\mathbf{I}$。球（用经纬卡 $\mathbf{c}_f(u,v)$）：$\mathbf{K}_f=\diagop(1/\rho,1/\rho)$、$\mathbf{M}_f=\diagop(\rho,\rho\cos u)$、$\mathbf{T}_f=[\,0,\ -\tfrac{1}{\rho}\tan u\,]$。代入纯滚动式 \cref{eq:gp-montana-roll}，相对曲率形式 $\mathbf{K}_f+\widetilde{\mathbf{K}}_o=\mathbf{K}_f+\mathbf{R}_\psi\mathbf{K}_o\mathbf{R}_\psi=\mathbf{K}_f=\tfrac1\rho\mathbf{I}$（平面 $\mathbf{K}_o=\mathbf{0}$），得
\begin{equation}
  \dot{\boldsymbol{\alpha}}_f=\mathbf{M}_f^{-1}\,\rho\begin{bmatrix}-\omega_y\\ \omega_x\end{bmatrix},\qquad
  \dot{\boldsymbol{\alpha}}_o=\mathbf{R}_\psi\,\rho\begin{bmatrix}-\omega_y\\ \omega_x\end{bmatrix}.
  \label{eq:gp-sphere-result}
\end{equation}
\textbf{读出物理}：平面上接触点的爬行速度 $\dot{\boldsymbol{\alpha}}_o$ 正比于\emph{球半径 $\rho$ 乘以滚动角速度}（$\rho\omega$ 正是球滚动的线速度），与平面在何处无关——印证“球滚一圈，地面轨迹长 $2\pi\rho$”的常识。这就是“平面上滚球，接触轨迹只由球半径决定”的精确来源。
\end{example}

\subsection{滚动抓取映射：接触运动学的下游}
当接触会滚动时，\cref{ch:mr-coopforce} 的抓取映射 $\mathbf{G}$ 不再是常数矩阵——它依赖接触坐标 $\boldsymbol{\eta}(t)$，须沿 Montana 方程（\cref{thm:gp-montana}）实时更新。

\begin{proposition}[滚动情形的抓取映射\cite{murray1994mathematical}]\label{prop:gp-rolling-grasp}
滚动接触下，抓取映射 $\mathbf{G}(\boldsymbol{\eta})$ 与手雅可比 $\mathbf{J}_h(\boldsymbol{\eta})$ 仍由“接触标架到物体/手指标架的旋量变换”给出（结构同固定接触，\cref{ch:grasp-closure}），但其分块为接触坐标的函数；接触坐标按 Montana 方程演化。固定（不滚动）接触是其中 $\dot{\boldsymbol{\eta}}=\mathbf{0}$、$\mathbf{G}$ 退化为常矩阵的特例。
\end{proposition}

\begin{insight}[滚动是“免重抓的重定位”]
Montana 方程的工程价值：它让机器人\emph{不松开物体}就能通过滚动改变物体在手内的姿态（in-hand reorientation）——把接触坐标 $\boldsymbol{\eta}$ 当作可控状态、滚动角速度当作输入，规划一条 $\boldsymbol{\eta}(t)$ 轨迹即得手内操作策略。这与 \cref{ch:arm-dualarm-planning} 在约束流形上规划、与 \cref{ch:mr-coopforce} 的内力调节，共同构成“持物状态下精细操作”的三块基石：滚动改\emph{接触}、约束流形改\emph{臂构型}、内力改\emph{夹持}。
\end{insight}

% ════════════════════════════════════════════════════════════════
\section{衔接：度量、构造与持物操作的全景}
\label{sec:gp-bridge}
% ════════════════════════════════════════════════════════════════

本章把 \cref{ch:grasp-closure} 的封闭判据延伸为度量、规划与非抓持操作。三条下游链路值得显式串起：

\begin{itemize}
  \item \textbf{到协同操作（\cref{ch:mr-coopforce}）}：本章的抓取力旋量集 $\mathcal{W}=\{\mathbf{G}\mathbf{f}\}$ 与 $\varepsilon$ 度量，正是 \cref{ch:mr-coopforce} 抓取矩阵 $\mathbf{G}$（\cref{eq:cm-grasp}，点接触 $\mathbf{G}_i=[\mathbf{I};\skewm{\mathbf{r}}_i]$，\cref{eq:cm-Gi}）的“质量评估”——多机/多指搬运在选定接触后，可用 $\varepsilon$ 判该抓取对外扰的鲁棒裕度；内力（$\ker\mathbf{G}$，见 \cref{eq:cm-decomp}）则是把接触力维持在摩擦锥内、从而维持本章封闭条件的旋钮。\cref{ch:mr-coopforce} 的接触模型分类（\cref{tab:cm-contact-models}）与本章 \cref{def:gp-gws} 的摩擦锥一一对应。
  \item \textbf{到闭链双臂规划（\cref{ch:arm-dualarm-planning}）}：刚性共持是“接触永不滚动/滑动”的极限，本章 Montana 方程退化为 $\dot{\boldsymbol{\eta}}=\mathbf{0}$（\cref{prop:gp-rolling-grasp}），抓取映射成常矩阵，闭链约束 $\mathbf{h}(\mathbf{q})=\mathbf{0}$ 即“相对位姿恒定”——这正是 \cref{ch:arm-dualarm-planning} 约束流形采样的出发点。本章的力封闭/内力分析回答“流形上某构型是否\emph{力}可行”，与那里“构型是否\emph{运动}可行”互补。
  \item \textbf{到双臂抓取对偶（\cref{ch:mr-multiarm}）}：本章度量可直接用于评估 \cref{ch:mr-multiarm} 中双臂抓取（\cref{eq:ma-grasp}）的鲁棒性。
\end{itemize}

% ════════════════════════════════════════════════════════════════
\section*{常见误解汇总}
% ════════════════════════════════════════════════════════════════
\begin{longtable}{p{0.40\linewidth} p{0.54\linewidth}}
\toprule
\multicolumn{1}{c}{误解} & \multicolumn{1}{c}{澄清} \\
\midrule
\endhead
“力封闭就一定夹得住，不会掉。” & 力封闭只说\emph{接触锥能生成任意力旋量}，不保证\emph{执行器真能提供}所需内力。\cite{lynch2017modern} 明确：电机出力不足或受限方向，力封闭物体仍可能掉（\cref{ex:gp-dynamic} 的反面）。\\
“抓取质量就是接触点数越多越好。” & 不是。质量看 $\mathcal{W}$ 的几何（$\varepsilon$、内接球），三个对称接触常优于五个挤在一侧的接触；冗余接触对 $\varepsilon$ 可能无贡献（\cref{thm:gp-steinitz} 的“可剔除冗余向量”）。\\
“$\varepsilon$ 度量是绝对的，可跨物体比较。” & 否。$\varepsilon$ 依赖力矩参考点、特征长度归一与 $L_1/L_\infty$ 约定（\cref{def:gp-gws} 的 note）。跨物体/跨约定比较前必须统一这三项。\\
“对踵条件对任意指数都成立。” & 对踵（\cref{thm:gp-antipodal}）是\emph{二指}充要；三指及以上须回 \cref{ch:grasp-closure} 的合成锥/LP 判据，无如此简洁的几何形式。\\
“无摩擦也能抓住球，只要接触够多。” & 不能。球是\emph{例外曲面}（\cref{prop:gp-contact-bounds} 的 note），所有法线过球心，绕心力矩无人可抗——无摩擦点接触无论多少都无法力封闭旋转面。\\
“推一个物体，它会沿推力方向直线走。” & 一般不会。底面分布摩擦经\emph{极限面}（\cref{def:gp-limit-surface}）合成力旋量，物体多半边走边转；转向由极限面 $+$ 推力作用线决定（\cref{ex:gp-push}）。\\
“滚动接触的抓取矩阵和固定接触一样是常数。” & 否。滚动时 $\mathbf{G}(\boldsymbol{\eta})$ 随接触坐标变，须沿 Montana 方程（\cref{thm:gp-montana}）更新；固定接触是 $\dot{\boldsymbol{\eta}}=\mathbf{0}$ 的特例。\\
\bottomrule
\end{longtable}

% ════════════════════════════════════════════════════════════════
\section*{本章小结}
% ════════════════════════════════════════════════════════════════
\begin{enumerate}
  \item \textbf{抓取质量}：把受限接触力经抓取映射得抓取力旋量集 $\mathcal{W}=\{\mathbf{G}\mathbf{f}\}$（凸多胞体，\cref{def:gp-gws}）。力封闭 $\iff\mathbf{0}\in\interior\mathcal{W}$。三个度量——最大内接球 $Q_{\text{ball}}=\min_k b_k$（\cref{prop:gp-ball-meaning}）、等价的 Ferrari--Canny $\varepsilon=\min_{\partial\mathcal{W}}\|\mathbf{w}\|$（\cref{def:gp-epsilon}，抗最坏外扰）、面向任务度量 $Q_{\text{task}}$（用任务力旋量集替球，\cref{def:gp-task}）——逐级特殊化。须处理力/力矩量纲、力矩参考点、$L_1/L_\infty$ 约定。
  \item \textbf{接触数界}：无摩擦点接触力封闭，Caratheodory 给下界 $p+1$、Steinitz 给上界 $2p$（平面 $4\!\sim\!6$、空间 $7\!\sim\!12$，\cref{prop:gp-contact-bounds}）；例外（旋转）面纯无摩擦永不可封闭。摩擦把空间手指数降到 $3$、平面到 $2$。
  \item \textbf{对踵与构造}：二指对踵定理（\cref{thm:gp-antipodal}）——连线落在两摩擦锥内 $\iff$ 力封闭；化为线性不等式（\cref{eq:gp-antipodal-ineq}）后可枚举平面多边形全部力封闭抓取、求最大独立接触区（\cref{algo:gp-construct}）。
  \item \textbf{非抓持操作}：统一用刚体动力学 $+$ 库仑摩擦 $+$ 接触模式枚举（\cref{algo:gp-contact-mode}），拟静态时退化为力旋量平衡（\cref{eq:gp-quasistatic}）。推动由摩擦极限面的法向法则决定转动中心（\cref{def:gp-limit-surface,ex:gp-push}）；动态抓取靠惯性力顶住物体（\cref{ex:gp-dynamic}）；meter-stick、峰嵌为拟静态/力封闭的生动副产物。
  \item \textbf{滚动接触}：由第一/二基本形式得几何参数 $(\mathbf{M},\mathbf{K},\mathbf{T})$（\cref{def:gp-first-form,def:gp-second-form,def:gp-gauss-frame}），Montana 方程（\cref{thm:gp-montana}）给接触坐标 $(\boldsymbol{\alpha}_f,\boldsymbol{\alpha}_o,\psi)$ 随相对旋量的演化；纯滚动时只受滚动角速度驱动（\cref{eq:gp-montana-roll}），球在平面滚的接触轨迹只由球半径定（\cref{ex:gp-sphere-plane}）。滚动使抓取映射成 $\mathbf{G}(\boldsymbol{\eta})$，是免重抓的手内重定位之本（\cref{prop:gp-rolling-grasp}）。
  \item \textbf{衔接}：度量评估 \cref{ch:mr-coopforce} 抓取矩阵 $\mathbf{G}$ 的鲁棒性、Montana 退化对接 \cref{ch:arm-dualarm-planning} 闭链约束流形、并可评 \cref{ch:mr-multiarm} 双臂抓取。
\end{enumerate}

\begin{finenote}
本章基于 Murray, Li \& Sastry《A Mathematical Introduction to Robotic Manipulation》第 5 章\cite{murray1994mathematical} 与 Lynch \& Park《Modern Robotics》第 12 章\cite{lynch2017modern} 的经典内容重述，记号统一到本书约定（抓取映射 $\mathbf{w}=\mathbf{G}\mathbf{f}$、点接触 $\mathbf{G}_i=[\mathbf{I};\skewm{\mathbf{r}}_i]$、$\SE(3)$ 旋量平移在前、右扰动）。对踵定理归 \citeauthor{nguyen1988}\cite{nguyen1988}，滚动接触方程归 \citeauthor{montana1988}\cite{montana1988}，$\varepsilon$ 度量归 \citeauthor{ferrari1992}\cite{ferrari1992}，极限面归 \citeauthor{goyal1991planar}\cite{goyal1991planar}。
\end{finenote}
